% Generated draft: requires mathematical and source review.
\documentclass[11pt]{article}
\usepackage[margin=25mm]{geometry}
\usepackage{fontspec}
\setmainfont{Latin Modern Roman}
\usepackage{amsmath,amssymb,mathtools,unicode-math}
\setmathfont{Latin Modern Math}
\usepackage{enumitem,needspace,xurl,hyperref}
\setlength{\parindent}{0pt}
\setlength{\parskip}{5pt}
\setlength{\emergencystretch}{3em}
\newcommand{\N}{\mathbb{N}}
\newcommand{\Z}{\mathbb{Z}}
\newcommand{\Q}{\mathbb{Q}}
\newcommand{\R}{\mathbb{R}}
\newcommand{\C}{\mathbb{C}}
\newcommand{\F}{\mathbb{F}}
\newcommand{\A}{\mathbb{A}}
\newcommand{\PP}{\mathbb{P}}
\newcommand{\E}{\mathbb{E}}
\newcommand{\eps}{\varepsilon}
\newcommand{\dd}{\,\mathrm{d}}
\begin{document}
\section*{UnsolvedMath ID: 800005}
\textbf{Source problem code:} catalogue entry
\subsection{Definitions and mathematical statement}
A semigroup is a nonempty set equipped with an associative binary operation. A variety of semigroups is a class of semigroups defined by identities in the semigroup language.

For context, under the usual equality-based interpretation, a semigroup \(S\) is residually finite if
\[
\forall a,b\in S,\quad a\ne b\Longrightarrow \exists\text{ a finite semigroup }F\ \exists\text{ a homomorphism }h:S\to F\text{ such that }h(a)\ne h(b).
\]
Under the corresponding usual local interpretation, a variety \(\mathcal V\) is locally residually finite if every finitely generated semigroup belonging to \(\mathcal V\) is residually finite. These definitions explain the conventional equality case only; the supplied record does not specify the predicate-dependent definitions intended in the full problem.

\textbf{Requested classification (S.~Kublanovsky).} Describe the varieties of semigroups that are locally residually finite ``for known predicates.'' The collection of predicates, their interpretations on semigroups, the corresponding residual finiteness conditions, and the desired form of the description are not specified in the supplied record. Consequently, the full requested classification cannot be expressed as an unambiguous quantified mathematical statement from this record alone.

\textbf{Supplied qualifications.} The source reports that the corresponding global descriptions were obtained in references [Kubl1] and [Kubl2], that the local description for the equality predicate was obtained in [Sapir], and that locally residually finite varieties of rings were described in [Kubl3]. These are supplied background claims, not independently verified results here. In particular, the equality case is presented as already described, so it must not silently be substituted for the entire unresolved target.
\subsection{Short English statement}
S.~Kublanovsky asks for a classification of locally residually finite semigroup varieties for the source's unspecified ``known predicates.'' The source reports prior descriptions in the global setting, for equality in the local setting, and for rings. The predicate-dependent definitions and scope must be recovered before a precise full formulation is possible. The registry labels this problem open; this draft does not independently verify that status.
\subsection{Sources}
\begin{itemize}
\item Local UnsolvedMath registry, permanent ID \(800005\), exact source code \texttt{catalogue entry}, title \emph{The Lyapin's Notebook --- Problem 5}. The supplied record belongs to \emph{AMR Open Problem Lists} and labels the problem \texttt{open}. Its upstream revision is \texttt{d371ee917ac9647fbb06cd6fc529deec6f3dda82}.
\item Bershadsky, Kublanovsky, and Mashevitzky, \emph{The Lyapin's Notebook: Unsolved Problems in Semigroup Theory} (2026), Problem 5, attributed in the statement to S.~Kublanovsky. Supplied source URL: \url{https://arxiv.org/abs/2604.04763}. The record reports access on 2026-10-16 and extraction by \texttt{arxiv\_latex}; these are supplied provenance details, not a claim of independent source access. The registry lists Bershadsky, Kublanovsky, and Mashevitzky in its proposer field and gives proposed year 2026.
\item References [Kubl1] and [Kubl2]: cited by the supplied statement for the global description. No titles, publication details, or individual URLs are supplied.
\item Reference [Sapir]: cited by the supplied statement for the local description for the equality predicate. No title, publication details, or individual URL is supplied.
\item Reference [Kubl3]: cited by the supplied statement for the description of locally residually finite varieties of rings. No title, publication details, or individual URL is supplied.
\end{itemize}
\end{document}
